\documentclass[10pt,a4paper]{amsart}
\usepackage[margin=19mm]{geometry}
\usepackage{amsmath,amssymb,mathtools}
\usepackage[scr=rsfs,cal=euler]{mathalfa}
\usepackage{xfrac}
\usepackage[unicode,hidelinks,bookmarks=false]{hyperref}
\newcommand{\ie}{i.e.\ }
\newcommand{\cf}{cf.\ }
\newcommand{\resp}{respectively\ }
\newcommand{\Subsection}{\subsection}
\providecommand{\nobreakdash}{\nobreak\hbox{-}}
\setlength{\emergencystretch}{2em}
\multlinegap0pt
\usepackage[shortlabels]{enumitem}
\usepackage{mathtools}
\usepackage{amssymb}
\usepackage{dsfont}
\usepackage{tikz}
\usepackage{float}
\newtheorem{lemma}{Lemma}
\newcommand\lcm{\mathrm{lcm}}
\title{The binary actions of\\ sporadic groups}
\author{Coen del Valle, Nick Gill}
\date{Source: arXiv:2609.02508v1 (CC BY 4.0)}
\begin{document}
\maketitle

\noindent\emph{Source and licence.} The binary actions of\\ sporadic groups. Version arXiv:2609.02508v1 (CC BY 4.0), distributed by arXiv under the Creative Commons Attribution 4.0 International licence, http://creativecommons.org/licenses/by/4.0/, as recorded in the arXiv licence metadata for this exact version and shown on the abstract page https://arxiv.org/abs/2609.02508v1. Full licence notice: \texttt{licenses/arxiv-2609.02508-CC-BY-4.0.txt}.\par\smallskip

\noindent\textit{Editorial scope.} Counted unit: the article's lemma lem:intrans (source0.tex lines 958--960). The article's complete original proof of it is delivered with it (source0.tex lines 961--971, one \texttt{proof} environment of 11 lines). No mathematics is written, repaired or re-derived: the statement and the proof are the article's own lines, in the article's own notation, and no retained line is substituted or re-flowed.  No further source-local context is needed: every cross-reference of every retained line resolves inside the two delivered spans.  Nothing was omitted from any retained span, so the package carries no omission record.  No retained line contains a citation, so the wrapper carries no bibliography and no bibliography entry.  No retained line contains a figure environment, an image-inclusion command or drawing code, so no image is delivered and the package declares no asset dependency.  Editorial formatting only, in addition: an AMS \texttt{amsart} preamble that restates the article's own declarations and macros used by the retained lines, namely the environments \texttt{lemma} on separate counters, together with the standard packages those lines need.  The retained environments print in this document as Lemma 1 in order of appearance, whereas the article prints the same items with the numbering of its own full document.  The complete native source of the article as distributed in the arXiv e-print for this exact version is bundled unchanged as \texttt{sources/209207/source0.tex}.
\begin{lemma}\label{lem:intrans}
    Let $g\in G\setminus Z(G)$ have prime order $p$, and let $1\ne z\in Z(G)$.     Suppose $H_1\leq \langle g,z\rangle$ and that the action of $G$ on $(G:\langle g,z\rangle)$ is binary.     Suppose that $|\langle g\rangle^G|\geq 3$ and that $g$ is conjugate to each of its non-trivial powers, but is not conjugate to $gt$ for any $t\in Z(G)\setminus \{1\}$. Fix integers $i$ and $j$ and suppose that $H_2\in\{\langle gz^i\rangle,\langle g,z^i\rangle\}^G$ and $H_3\in \{\langle gz^j\rangle,\langle g,z^j\rangle\}^G$ are such that $Z(G)H_1, Z(G)H_2,$ and $Z(G)H_3$ are distinct. Then $H_1\cap (H_2\cdot H_3)\subseteq H_1\cap((H_1\cap H_2)\cdot H_3)$.
\end{lemma}

\begin{proof}
     Since $Z(G)H_1,Z(G) H_2,$ and $Z(G)H_3$ are distinct, we may find conjugates $K_1=\langle g,z\rangle\geq H_1$, $K_2=\langle g_2,z\rangle\geq H_2$, and $K_3=\langle g_3,z\rangle\geq H_3$ of $\langle g,z\rangle$ where no two of the cyclic subgroups $\langle g_i\rangle$ are equal (using that $|\langle g\rangle^G|\geq 3$). Therefore, since the action of $G$ on $\langle g,z\rangle $ is binary, \begin{align*}
            H_1\cap (H_2\cdot H_3)&\subseteq K_1\cap (K_2\cdot K_3)\\&\subseteq K_1\cap ((K_1\cap K_2)\cdot K_3)\\&=K_1\cap (\langle z\rangle\cdot K_3)\\&\subseteq\langle z\rangle.
        \end{align*}
    
     We will assume that $H_2=\langle g_2z^i\rangle$ and $H_3=\langle g_3,z^j\rangle$ for some $g_2\ne g_3\in g^G$; the other cases are very similar. Suppose $1\ne z_0\in H_1\cap (H_2\cdot H_3)\subseteq\langle z\rangle$. Then $z_0=(g_2^mz^{im})(g_3^nz^{jk})$ where $m,n,k$ are integers. Then, $g_2^m\in g_3^{-n}Z(G)$, so since $g$ is not conjugate to $gt$ for any $t\in Z(G)\setminus{1}$ we deduce that $g_2^m=g_3^{-n}$. But $\langle g_2\rangle\cap\langle g_3\rangle=1$, so $p\mid m$, hence $z_0=z^{im+jk}\in \langle z^{\gcd(pi,j)}\rangle\cap H_1$. Thus, if $z^s$ generates $H_1\cap Z(G)$, then $z_0\in \langle z^{\lcm(\gcd(pi,j),s)}\rangle$.
    Finally, \begin{align*}
             H_1\cap ((H_1\cap H_2)\cdot H_3)&= H_1\cap (\langle z^{\lcm(pi,s)}\rangle\cdot H_3)\\&=H_1\cap \langle g_3,z^{\gcd(\lcm(pi,s),j)}\rangle\\&= \langle z^{\lcm(\gcd(\lcm(pi,s),j),s)}\rangle.
        \end{align*}
    Now, $\lcm(\gcd(pi,j),s)=\gcd(\lcm(pi,s),\lcm(j,s))=\lcm(\gcd(\lcm(pi,s),j),s)$, from the distributive properties of $\lcm$ and $\gcd$, hence the result.
   \end{proof}

\end{document}
